\chapter{Getting Started}
\label{ch:getting-started}

\section{Installing Calango}

Calango ships as a native installer for each desktop platform:

\begin{description}
  \item[macOS] A drag-and-drop \file{.dmg}. Mount it and drag
    \file{calango.app} onto the \file{Applications} shortcut.
  \item[Linux] A Debian package. Install it with \opt{apt} so dependencies
    resolve automatically:
    \begin{lstlisting}
sudo apt install ./calango_26.7.26_amd64.deb
    \end{lstlisting}
    The package registers a desktop launcher, an application icon, and a
    MIME type for \file{.calproj} files, so double-clicking a project in
    your file manager opens the workspace.
\end{description}

Building from source is covered in the companion \emph{Calango Packaging
Guide} (\file{docs/packaging.pdf}), which also documents how the installers
are produced.

\section{The Python environment}
\label{sec:python-env}

Calango embeds a CPython interpreter and drives ASE through it. Almost
everything that touches chemistry --- reading a CIF, building a slab,
running a calculation --- goes through that interpreter, so pointing
Calango at a Python that has ASE installed is the single most important
piece of setup.

\subsection{How the interpreter is chosen}

An embedded interpreter does not inherit an activated virtual environment
the way a shell does, so Calango resolves one explicitly, in this order:

\begin{enumerate}
  \item \opt{\$CALANGO\_PYTHON} --- an explicit path to a Python
    executable. Highest priority; use this to override everything else.
  \item \opt{\$VIRTUAL\_ENV} --- an activated virtual environment, if
    Calango was launched from that shell.
  \item A Python bundled inside the installed application
    (\file{Contents/Resources/python} on macOS,
    \file{\textless prefix\textgreater/lib/calango/python} on Linux), when the installer
    was built with one.
  \item The interpreter the build was configured against.
\end{enumerate}

\begin{calwarn}
Calango needs \texttt{ase} and \texttt{numpy} at minimum. Individual
features add their own requirements: \texttt{spglib} (symmetry, Raman
modes), \texttt{paramiko} (remote access), \texttt{pillow} /
\texttt{imageio} / \texttt{imageio-ffmpeg} (animation export),
\texttt{mace-torch} (ML potentials), \texttt{gpaw} (DFT bands and PDOS). A complete environment:
\begin{lstlisting}
python3 -m venv .venv
.venv/bin/pip install ase numpy spglib paramiko pillow imageio \
                      imageio-ffmpeg
\end{lstlisting}
\end{calwarn}

\subsection{Checking the environment}

Before opening any structure, confirm the interpreter is the one you
expect:

\begin{lstlisting}
calango --probe-python
\end{lstlisting}

This runs headlessly and prints the resolved interpreter, the Python
version and the ASE version, exiting with status 0 only when ASE imports
successfully. It is the first thing to run when structure loading fails.

If ASE cannot be imported at launch, Calango still starts --- so you can
adjust settings --- but structure I/O and job features are disabled, and a
warning explains how to point at a working interpreter.

\begin{caltip}
Simulation jobs do not have to use the embedded interpreter. The
calculator, phonon and band-structure dialogs each carry an
\ui{Execution Environment} selector that routes the job to any Conda
environment or interpreter path (Section~\ref{sec:exec-env}). This is how
you keep, say, a GPAW environment separate from the one Calango itself
uses.
\end{caltip}

\subsection{The environment file}
\label{sec:envfile}

At launch Calango reads an environment file --- \file{\textasciitilde/.env}
by default --- and exports \opt{MP\_API\_KEY}, the Materials Project API
key, into the process environment. A key already present in your shell
takes precedence at startup.

Point Calango at a different file, or re-read it on demand, in
\menu{Edit}\menusep\menu{Preferences}. The dialog reports whether the file
exists and whether it actually defines a key. The same controls are
mirrored in the \ui{Materials Project} tab of the database browser.

\section{First launch}

Calango opens with an empty workspace. Three quick ways to get something
on screen:

\begin{enumerate}
  \item \menu{File}\menusep\menu{Open}\menusep\menu{Structure} and pick any
    structure file (Section~\ref{sec:formats} lists the formats).
  \item \menu{Build}\menusep\menu{From Database} and load one of the
    bundled presets --- diamond, bulk and monolayer MoS$_2$, graphene,
    benzene, naphthalene, coronene.
  \item Pass a file on the command line: \lstinline|calango structure.cif|.
    Each file argument opens in its own tab; a \file{.calproj} argument
    restores a whole saved session.
\end{enumerate}
